% \label{sec:3.3}

We present a Knowledge Path Retriever to obtain context-relevant knowledge paths from all possible knowledge paths motivated by \cite{lewis2020retrieval,kang2022knowledge}.
Given a dialog history $\TokenSeq$ and knowledge paths $\PathSet$, the knowledge path retriever computes a relevance score between the dialog history and each path and returns the most relevant $k$ knowledge paths.
Similar to the knowledge path embedding $\hb_p = \enc(\zb_p)$ in \eqref{eq:kpembed}, 
we feed the token sequences of dialog history $\TokenSeq$ into a Transformer encoder $\enc(\cdot)$ and compute the context vector $\hx$, which is \texttt{[CLS]} output vector, \ie, $\hb_x = \enc(\xb)$.
With the context embedding $\hb_x$ and the knowledge path embedding $\hb_p$, 
the context-relevance score between the dialog history $\xb$ and path $p$ is computed as:
\begin{equation}
s ( \xb, p ) = \mathbf{\sigma\left( \MLPx (h_{\TokenSeq})^\top MLP_p\left(h_p \right)\right)},
\label{eq:kpscore}
\end{equation}
where $\sigma$ is the sigmoid function, $\MLPx$ and $\MLPp$ are single-layer feed-forward networks to project embedding vectors onto a common space.\footnote{Note that we only fine-tune the feed-forward networks and freeze the dialog history encoder for the retrieval.}
Finally, based on the relevancy scores $s (\xb, p)$, the $k$ most context-relevant knowledge paths are retrieved. 
This Knowledge Path Retriever can be viewed as a path-level hard attention layer, which suppresses irrelevant paths and lets the model more focus on informative paths. More discussion with experimental results is in Section \ref{sec:analysis}.
We then retrieve informative paths $\hat{\PathSet}$ by selecting paths with the top $k$ relevance scores.